\subsection{Scope and caveats}\label{subsec: caveats}

\paragraph{Calibration identity versus existence of a calibrated model.} Propositions \ref{prop: calibration} and \ref{prop: main} are statements about a model that exists: \emph{if} (\ref{eq: threeFactorModel}) has a sufficiently regular solution, then it is calibrated if and only if its local volatility satisfies (\ref{eq: mainFormula}). Substituting (\ref{eq: mainFormula}) into (\ref{eq: threeFactorModel}) gives an SDE which is nonlinear in the sense of McKean, whose coefficient depends on the law of the solution through a conditional expectation given $S_t$. Neither existence nor uniqueness of a solution of this equation follows from the identity. The existence results above cover the stated Gaussian subclass and a regularized equation, not arbitrary unregularized smiles. Even for local-stochastic volatility models with deterministic rates, existence is known only in special cases: perturbative and local in time \cite{abergelTachet}, for regime-switching volatility under a condition on the range of the volatility factor \cite{jourdainZhou}, and for stationary solutions of a time-homogeneous version \cite{lackerShkolnikovZhang}. The more recent results of \cite{djete,mustapha} also concern different LSV structures and do not directly establish existence for the present general-smile stochastic-rates equation. The particle algorithm is therefore a numerical construction; its convergence as $N\to\infty$ is not proved in this paper.

\paragraph{Non-negativity of the local variance.}
The right hand side of (\ref{eq: mainFormula}) need not be positive.
The explicit Gaussian sufficient condition and the general necessary
log-contract bound in Section \ref{subsec: existence} have different scopes:
failure of a particular Gaussian construction is not a proof of nonexistence
among all spot-dependent models. In the numerical algorithm, report where
and by how much the variance floor binds. Such activation may reflect
estimation, interpolation or time-discretization error as well as incompatible
inputs; a negative trial iterate alone does not distinguish these cases.

\paragraph{Regularity.}
The hypotheses (H1)--(H3) distinguish local flow regularity and Malliavin
integrability from the density conditions needed for a classical strike-wise
formula. A global bound on $S\partial_S\sigma$ and a nonsingular correlation
matrix are not needed for smooth-payoff duality. In the algorithm, however,
noise in the estimated derivative $\partial_S\sigma_N$ enters the auxiliary
processes multiplicatively. Relaxing analytical assumptions does not remove
this numerical sensitivity.

\paragraph{What has and has not been tested.} The Monte Carlo evidence of Section \ref{subsec: checks} concerns the identities behind (\ref{eq: mainFormula}) for a \emph{given} smooth local volatility, and the simulation scheme. The full particle calibration has been run once, on CME EUR/USD option settlements of 12 November 2026: expiries from December 2026 to December 2027, the surface being extended beyond that at constant implied-volatility smile to five years, which is synthetic. The calibration is run under the domestic risk-neutral measure with discount weights, 200,000 particles and 100 time steps per year, and two sets of rate parameters which are illustrative and not calibrated to interest rate options. The options are then repriced with an independent simulation. Over the quoted expiries the repricing error attributable to the stochastic rates is at most 0.005 volatility points with the correction (\ref{eq: mainFormula}), against 0.19 and 0.47 for the two parameter sets with the Dupire local volatility alone. With the more volatile set, the calibrated variance reaches its floor within one standard deviation of the money from about three years, on the synthetic part of the surface, and the five-year at-the-money error is 0.54. The report and the scripts are in the repository \texttt{github.com/ceikit/sir-lsv-malliavin-particle}, file \texttt{calibration/RESULTS.md}.

A second test is synthetic and uses options on the spot. A non-flat local volatility, whose level and skew vary in time, generates the target Dupire variance on a two-year horizon. The volatility is then recovered from that target by three methods: a three-dimensional forward equation for the discounted density with the conventional tail integrals, the five-field system of Section \ref{subsec: momentClosure}, and the particle method with the sensitivities (\ref{eq: sensitivities}). On their finest settings the three methods recover the generating volatility with a root-mean-square error between 0.002 and 0.003 volatility points, and the calibrated models reprice options against the generating model within 0.004. With 20,000 particles the particle method has an error of 0.004 after about three seconds; the two PDE solvers take 69 and 189 seconds on their finest grids. This is one set of parameters and one implementation of each method, and at this accuracy the error of the reference target is not negligible, so the methods cannot be ranked. Details are in \texttt{benchmarks/RESULTS.md}.

Several things remain untested. The CME options are options on futures and were treated as options on the forward to their expiry. The bias due to the time step and to the kernel bandwidth has not been quantified systematically. No market data beyond December 2027 was used, and no comparison has been made with other particle methods such as \cite{cozma,ogetbilHientzsch}.

\paragraph{Hull-White.} Proposition \ref{prop: main} uses two properties of the Hull-White model: the Clark-Ocone integrands (\ref{eq: kernels}) are deterministic and factorise, and the bond volatilities are deterministic, so that risk-neutral and forward measures differ by deterministic drifts.

\paragraph{Related literature.} The use of Malliavin integration by parts to remove the density from the stochastic-rates correction is not new. Bloch and Nakashima \cite{blochNakashima}, who credit \cite{benhamouRivoiraGruz}, treat the FX local volatility with stochastic domestic and foreign rates in this way, display a correction of the same domestic-minus-foreign form as (\ref{eq: mainFormula}), and calibrate by a fixed-point iteration. A comparison at the level of the equations (the note on \cite{blochNakashima} in \texttt{docs/reviews/} of the repository) found that their implemented two-rate method freezes the drift locally and evaluates the volatility at a fixed strike, and found no equivalent of the four forward equations (\ref{eq: sensitivities}), in particular of the terms produced by the foreign change of measure. What this section may add is therefore narrow: an exact finite-dimensional realisation of the correction for the coupled Gaussian-rates model, and the calibration equivalence of Proposition \ref{prop: main}. The paper \cite{benhamouRivoiraGruz} itself has not been read in full, and no claim of priority is made.

\subsection{Outlook}
\paragraph{Existence.} In the system (\ref{eq: closureFive}), the equations of the two moments of $\ell^d$ and $\ell^f$ can be rewritten without second derivatives of the unknown variance $\sigma^2$ (a note in \texttt{paper/notes/} of the repository). Together with the known-denominator formulation of Theorem \ref{thm: closureNonlinear}, this is a possible starting point for a perturbative existence argument around the flat-smile solution of Proposition \ref{prop: gaussianExistence}. The estimates such an argument needs have not been obtained.

\paragraph{Gaussian multi-factor rates.} If a rate is a sum of Hull-White factors, the kernel (\ref{eq: kernels}) is a sum of exponentials, one per factor, and the matrix system (\ref{eq: matrixSystem}) extends with one Cameron-Martin direction per factor. The reduction to scalar processes has not been carried out for that case. Factors with different mean reversions cannot be aggregated into one process, and the cross sensitivities between each domestic and each foreign factor have to be tracked separately, so the number of auxiliary processes is not simply two per factor.

\paragraph{General one-factor short rates.} For rates $dr^i_t=\mu^i(t,r^i_t)dt+\sigma^i_r(t,r^i_t)dB^i_t$ the statement that replaces the factorisation $V_t/V_s$ is $\MD_sX_t=J_tJ_s^{-1}\mathcal{A}(s,X_s)\mathcal{C}^\top e$, with $J$ the $3\times3$ first variation matrix of $X=(S,r^d,r^f)$, which is not triangular because of the quanto drift. Two further difficulties appear: the Clark-Ocone integrand is stochastic and depends on $t$, and the forward measures differ from the risk-neutral ones by state-dependent drifts, which produce additional terms of the type of $H_t$.

\paragraph{Stochastic volatility.} With a stochastic volatility factor $a_t$ as in Section 1, the left hand side of (\ref{eq: mainFormula}) becomes $\sigma(t,K)^2\,\mathbb{E}^{d^t}[a_t^2|S_t=K]$ and the sensitivities acquire the tangent process of $a_t$ when $a_t$ is correlated with the rates. This is not worked out here.
